\begin{textsample}{POS Dim 1 – sc – Score 35.00 – sc/sc\_inf\_e\_ef\_33.txt}  \label{ex:f1_pos_020}
5.1.3 \textbf{Encaminhamentos} para as indicações \textbf{metodológicas} Como mencionado, ao término do estudo dos Quadros 1 a 9, você, professor(a), pode questionar sobre a possibilidade de se contemplar todos os objetos de conhecimento e todas as habilidades definidas para cada ano, pela BNCC.
A resposta vai \textbf{depender} do modo de organização do ensino.
Se o(a) professor(a) for abordar objeto por objeto e suas correspondentes habilidades, separadamente, provavelmente não será possível concluir nem o mínimo definido pela BNCC, para cada ano.
Além disso, cabe outro questionamento : os objetos e as habilidades, anteriormente apresentados, serão desenvolvidos em sala de aula em nível empírico ou \textbf{teórico}?
No Estado de Santa Catarina, \textbf{historicamente}, a opção foi e continua sendo pela apropriação dos conhecimentos científicos e pelo desenvolvimento do pensamento \textbf{teórico}.
Nesse sentido, as habilidades anteriormente apresentadas consistem em manifestações particulares da relação universal, essencial e nuclear dos conceitos e dos sistemas \textbf{conceituais}.
O desenvolvimento de habilidade por habilidade, em sala de aula, de forma fragmentada, resultará no desenvolvimento do pensamento empírico.
Por outro \textbf{lado}, se forem considerados os princípios didáticos decorrentes dos \textbf{fundamentos} \textbf{teóricos} da Proposta Curricular do Estado de Santa Catarina, não só é possível dar conta do estabelecido pela BNCC, mas também ir além.
De acordo com os \textbf{fundamentos} e os \textbf{desdobramentos} da Teoria Histórico-Cultural para o ensino, os objetos de conhecimento não são abordados de forma fragmentada, mas em conexão com outros do sistema \textbf{conceitual} no qual se insere, a partir de sua relação geneticamente inicial, universal e essencial.
Nesse movimento \textbf{teórico}, as habilidades anteriormente apresentadas consistem em manifestações particulares, decorrentes da gênese \textbf{que} precisa ser \textbf{revelada} e modelada durante o processo de ensino e aprendizagem.
Nessa \textbf{direção}, a seguir, nas indicações \textbf{metodológicas}, apresenta -se \textbf{um} exemplo de caráter geral \textbf{que} integra alguns objetos e habilidades previstas para cada ano do Ensino Fundamental, à luz dos \textbf{fundamentos} e dos \textbf{desdobramentos} da Teoria Histórico-Cultural, com a finalidade de orientar o(a) professor(a) na elaboração de suas próprias situações desencadeadoras de aprendizagem para o ensino de Matemática.
Trata -se de \textbf{uma} Situação Desencadeadora de Aprendizagem \textbf{que} pode ser desenvolvida em qualquer ano do Ensino Desenvolvimental, pois, como é geral, pode ser resolvida tanto por meio de conceitos mais básicos, previstos para os primeiros anos de escolarização, como por meio de conceitos mais abstratos, correspondentes aos últimos anos do Ensino Fundamental. 5.1.4 Indicações \textbf{metodológicas} Avançar no desenvolvimento científico e tecnológico à margem do pensamento matemático no âmbito \textbf{teórico} é \textbf{uma} \textbf{tarefa} difícil.
Assim sendo, é recomendável elaborar situações significativas \textbf{que} \textbf{desencadeiem} a apropriação pelos estudantes do \textbf{que} de mais \textbf{atual} a \textbf{humanidade} produziu em \textbf{termos} de conceitos matemáticos e recursos tecnológicos, de modo indissociável.
Esta unidade constitui \textbf{um} dos \textbf{alicerces} indispensáveis ao processo de formação integral dos \textbf{sujeitos}, \textbf{que} passa pela necessidade de se aprofundar os conceitos matemáticos em nível científico e pela transformação das escolas em ambientes sustentáveis e tecnologicamente \textbf{atuais}. \textbf{Contudo}, é importante enfatizar \textbf{que} a existência de recursos tecnológicos de última geração e \textbf{uma} Situação Desencadeadora de Aprendizagem bem elaborada não garantem, por si só, a apropriação de conceitos científicos e o desenvolvimento \textbf{teórico} por parte dos estudantes.
Isso \textbf{depende} da lógica \textbf{que} o professor adota para orientar as reflexões individuais e coletivas.
As Situações Desencadeadoras de Aprendizagem podem ser desenvolvidas por meio de recursos tecnológicos de ponta empírica ou teoricamente, do ponto de vista do conhecimento matemático – vai \textbf{depender} da lógica \textbf{que} sustenta o movimento \textbf{conceitual} a ser conduzido pelo professor ( \textbf{SANTA} CATARINA, 2014 ).
No \textbf{que} \textbf{diz} respeito à diversidade humana e à formação integral do ser humano, dentre outros aspectos, a Proposta Curricular de Santa Catarina ( 2014, p. 25 ) \textbf{diz} \textbf{que} “ [... ] quanto mais integral a formação dos \textbf{sujeitos}, maiores são as possibilidades de criação e transformação da sociedade ”.
Com esse propósito, a Situação Desencadeadora de Aprendizagem a seguir, apresentada na forma de \textbf{uma} história, pode ser desenvolvida em decorrência das reflexões sobre a diversidade de “ arranjos familiares \textbf{hoje} possíveis ” ( \textbf{SANTA} CATARINA, 2014, p. 59 ).
Situação Desencadeadora de Aprendizagem : os preparativos para o dia da família na escola de João e Maria Para organizar o dia da família na escola, as \textbf{tarefas} foram distribuídas por turma.
A turma de João e Maria ficou encarregada de determinar a quantidade de fita decorativa necessária para fixar a(s) toalha(s) ao redor da(s) mesa(s ).
A única informação \textbf{que} a turma recebeu foi \textbf{que} deveriam utilizar, como unidade padrão, \textbf{uma} das medidas da(s) mesa(s ).
A data da festa já estava se \textbf{aproximando} e a turma ainda não sabia como proceder.
Diante desse contexto, como podemos auxiliar a turma de João e Maria a determinar a quantidade necessária de fita para fixar a(s) toalha(s) ao redor da(s) mesa(s)?
Para a contação da história na presença de estudantes surdos ou surdocegos, faz -se necessária a realização de adequações a fim de garantir a interação entre todos os estudantes.
Além disso, é importante destacar \textbf{que} a inclusão de todos os estudantes nas reflexões individuais e coletivas é condição primordial para o desenvolvimento de \textbf{uma} Situação Desencadeadora de Aprendizagem em sala de aula.
Isso passa pela inclusão digital, inclusão dos estudantes com necessidades educativas específicas, entre outras.
Afinal, é \textbf{preciso} garantir o envolvimento e a participação ativa de todos nas reflexões individuais e coletivas, em \textbf{um} contexto democrático, humanizador e sustentável.
O currículo catarinense ressalta : > [... ] a necessidade do respeito à diversidade humana em todas as suas múltiplas dimensões. [... ] para \textbf{que} todos os \textbf{sujeitos} \textbf{que} integram as comunidades, escolares ou não, tenham respeitadas sua \textbf{dignidade} e direito de opção, seja ela voluntária ou ditada pela própria natureza humana.
Outro aspecto a ser considerado é a \textbf{percepção} de \textbf{totalidade} de mundo, \textbf{homem}, sociedade, necessária aos professores [... ]. ( \textbf{SANTA} CATARINA, 2014, p. 172 ).
A temática da Situação Desencadeadora de Aprendizagem foi definida a partir de \textbf{um} evento \textbf{que} já faz parte da cultura escolar catarinense : o dia da família na escola.
Trata -se de \textbf{uma} situação em \textbf{que} os estudantes precisam pensar como se faz para calcular a quantidade de fita necessária para prender as toalhas nas mesas.
Envolve, portanto, a relação entre a grandeza comprimento ( perímetro da superfície das mesas e \textbf{uma} de suas partes a ser considerada como unidade de medida ) e a grandeza discreta ( quantidade de mesas ).
Consiste em \textbf{uma} situação de caráter geral e, por \textbf{conseguinte}, pode ser desenvolvida em qualquer \textbf{um} dos anos escolares, de modo \textbf{que} possibilite o acompanhamento da progressão de conhecimento, tal como está previsto na BNCC ( BRASIL, 2017 ).
Trata -se de \textbf{um} esforço por exemplificar algumas possibilidades de articulação da BNCC ( BRASIL, 2017 ) com a Proposta Curricular de Santa Catarina ( 2014 ).
O ponto de partida na sua elaboração foi o núcleo do conceito \textbf{que} dá origem aos números naturais, racionais e irracionais : a relação universal de multiplicidade e divisibilidade, \textbf{logo}, o número c pode ser natural, racional ou irracional.
Essa relação universal dá origem a \textbf{um} sistema de conceitos \textbf{que} perpassa todos os anos do Ensino Fundamental e pode ser introduzido a partir da Situação Desencadeadora de Aprendizagem.
No Quadro 10 a seguir, apresentam -se as orientações gerais ( primeira coluna ), particulares ( segunda coluna ) e singulares ( terceira coluna ), para o desenvolvimento de \textbf{uma} Situação Desencadeadora de Aprendizagem em nível teórico/científico.
Todas as ações e as operações apresentadas no quadro a seguir são desenvolvidas por meio dos mais variados recursos tecnológicos, a \textbf{depender} das condições objetivas da escola e o perfil dos estudantes subsidiados pela diferenciação curricular e desenho universal de Aprendizagem ( DUA ).

% matched lemmas: alicerce, aproximar, atual, conceitual, conseguinte, contudo, depender, desdobramento, desencadear, dignidade, direção, dizer, encaminhamento, fundamento, historicamente, hoje, homem, humanidade, lado, logo, metodológico, percepção, preciso, que, revelar, santo, sujeito, tarefa, termos, teórico, totalidade, um
\end{textsample}
